\chapter{Conclusiones}

Este capítulo ofrece una conclusión provisional pero utilizable: cierra lo que ya está sustentado por el diseño, la implementación y los artefactos históricos, y deja condicionadas las afirmaciones que dependen de la campaña final. No se presenta el protocolo aprobado como si fuese un resultado ejecutado.

\section{Problema abordado y sistema construido}

El trabajo aborda una pregunta concreta: cómo formular y evaluar, de manera auditable, un defensor offline que decida flujo a flujo entre \texttt{PERMIT} y \texttt{BLOCK} cuando los errores tienen costes distintos. La dificultad no reside solo en alcanzar una métrica alta. Un NIDS experimental debe evitar fugas, distinguir rendimiento intradistribución de generalización, conservar el preprocesamiento exacto y limitar su alcance frente a un benchmark conocido por su heterogeneidad y duplicados.

La respuesta construida comprende un cargador de CICIDS2017, un contrato canónico de 76 características y 76 indicadores de presencia, exclusión de identificadores susceptibles de fuga, escalado ajustado solo con entrenamiento, un entorno Gymnasium, una recompensa asimétrica, un agente QRDQN y una línea base de Random Forest. El modelo y sus evaluaciones producen artefactos trazables, y la Fase~2 reutiliza el mismo contrato para inferencia offline sobre flujos de laboratorio. No se ha construido un cortafuegos autónomo ni un sistema de producción.

\section{Contribución metodológica}

La contribución más sólida es el diseño experimental reproducible. El perfil \texttt{main-v1} congela la configuración científica; \texttt{split\_seed} y \texttt{model\_seed} separan la partición de la aleatoriedad del aprendizaje; la caché conserva datos canónicos sin escalar; cada ejecución ajusta su propio escalador solo con entrenamiento; y el esquema de artefactos registra configuración, métricas, entorno, tiempos, monitorización, modelo, preprocesamiento, manifiesto y sumas de comprobación. La campaña final organiza estas piezas en contrastes que responden a preguntas distintas, no en una colección de ejecuciones aisladas.

También es una contribución declarar con precisión la naturaleza del problema. Con $\gamma=0$, QRDQN actúa como un bandido contextual sensible al coste. Esta elección permite estudiar una política de decisión y su distribución de retornos, pero no atribuir al agente planificación temporal ni consecuencias activas sobre una red. La honestidad de esta interpretación mejora la comparabilidad con modelos supervisados y evita exagerar el alcance de RL.

\section{Qué respalda la evidencia disponible}

Los artefactos históricos respaldan tres conclusiones provisionales. Primero, el pipeline puede aprender la tarea binaria bajo una partición aleatoria de CICIDS2017 y conservar todo lo necesario para auditar esa evaluación. Segundo, los controles históricos con etiquetas barajadas, validación directa y análisis de duplicados justifican que el rendimiento aleatorio se interprete como comprobación interna y no como prueba suficiente de generalización. Tercero, Random Forest es una referencia fuerte: puede superar a QRDQN en distribución y, por tanto, impide defender que RL sea preferible por definición.

La inferencia histórica de Fase~2 respalda una conclusión de ingeniería limitada: el modelo y el preprocesamiento persistidos pudieron aplicarse a una captura concreta externa al CSV de entrenamiento. No respalda robustez frente a producción, adversarios independientes ni otros entornos.

\section{Qué no respalda todavía}

La evidencia disponible no permite afirmar estabilidad entre semillas, eficiencia de datos, robustez del perfil completo ante la separación por días, comportamiento consistente en los cuatro escenarios retenidos, superioridad de QRDQN frente a RF ni transferencia final de la MAIN fresca. Tampoco permite elegir un umbral operativo, estimar impacto real de bloqueos o declarar preparación para despliegue. Los intervalos \emph{bootstrap} históricos miden precisión sobre un test fijo, no varianza de entrenamiento.

Estas ausencias no son meras tareas editoriales. Definen el límite actual de la conclusión científica. Por ese motivo el resumen no contiene una cifra final y los capítulos de resultados y discusión conservan estructuras de inserción, no valores provisionales.

\section{Por qué la partición aleatoria no basta}

Una partición aleatoria mezcla filas procedentes de días y escenarios comunes y puede situar duplicados exactos a ambos lados de la frontera. Sirve para comprobar que el pipeline aprende y para comparar modelos bajo una distribución controlada. No exige enfrentarse a una captura distinta. La partición completa por días rompe parte de esa continuidad; los \emph{holdouts} aíslan escenarios específicos; la sensibilidad a semilla separa un resultado puntual de una tendencia local; y la escalera de tamaños muestra si el comportamiento depende del volumen y del presupuesto de entrenamiento. Solo la lectura conjunta de estos contrastes puede sostener una conclusión de generalización proporcionada a la evidencia.

\section{Papel de Random Forest}

Random Forest cumple una función interpretativa central. En una tarea tabular de un paso, obliga a demostrar que QRDQN aporta algo más que un clasificador complejo aplicado a buenas características. La comparación final será emparejada por partición, tamaño y escenario. Si ambos modelos siguen el mismo patrón, la explicación más probable residirá en el conjunto de datos o en el protocolo. Si difieren de forma reproducible, podrá discutirse una propiedad específica del aprendizaje o de la política sensible al coste. En ningún caso una sola partición justificará una afirmación universal.

\section{Aprendizajes sobre RL con un conjunto de datos como entorno}

El trabajo muestra que un conjunto etiquetado puede envolverse como entorno RL de forma reproducible, pero que esta transformación no crea por sí misma dinámica secuencial. La recompensa obliga a expresar preferencias entre falsos positivos y falsos negativos; el agente permite estudiar una política distribucional; y los artefactos de RL hacen observable el entrenamiento. Aun así, el valor científico depende de la calidad de la partición, la ausencia de fugas y la comparación supervisada, igual que en clasificación. El principal aprendizaje es que la formulación RL debe evaluarse con más controles, no con menos.

\section{Cuestiones reservadas a la campaña final}

La campaña está diseñada para resolver cinco cuestiones pendientes: si una MAIN fresca reproduce el comportamiento intradistribución; si este persiste con separación completa por días; cómo cambia con el tamaño de entrenamiento; cuánto varía entre semillas de modelo; y qué escenarios retenidos provocan degradación. Los RF emparejados determinarán si los patrones son específicos de QRDQN, mientras que los controles directos, \emph{bootstrap}, duplicados y etiquetas barajadas reforzarán la validez interna. La Fase~2 fresca cerrará la cadena con el modelo final.

\evidencenote{La versión de entrega deberá sustituir esta nota por una síntesis breve, derivada de los artefactos finales, que responda a las cinco cuestiones anteriores y mantenga explícitas las amenazas que no queden resueltas.}

\section{Cierre provisional}

Antes de disponer de la campaña final, la conclusión defendible no es que QRDQN haya resuelto la detección de intrusiones ni que supere a un clasificador supervisado. Es que el proyecto ha convertido una propuesta de defensor basado en flujos en un método verificable: define qué observa, qué decide, qué coste optimiza, cómo evita fugas, contra qué se compara y qué evidencia necesita para hablar de generalización. Esa estructura permite que los resultados finales, sean favorables o adversos, puedan incorporarse sin cambiar la pregunta científica ni ocultar sus límites.

Como trabajo futuro posterior a la campaña quedan la validación con capturas independientes de varios entornos, más repeticiones y familias de modelo, tratamiento de tráfico desconocido y una evaluación de seguridad específica antes de considerar acciones reales. Ninguna de estas extensiones puede deducirse automáticamente de CICIDS2017 ni de un laboratorio cerrado.
